\section{Experiment Results}
\label{sec:results}

\subsection{Experimental Setup}

\paragraph{Dataset.}
We study 10 peer-reviewed modeling papers covering cholera~\cite{Fung2014Cholera},
COVID-19~\cite{TuiteE497}, dengue~\cite{dengue}, Ebola~\cite{Abbas2024Ebola},
influenza~\cite{Corbella2017MonitoringAP}, HIV~\cite{Espitia2022},
malaria~\cite{Akowe2025Malaria}, measles~\cite{Pokharel2024Measles},
tuberculosis~\cite{tuberculosis}, and Zika~\cite{Lee2016Zika}.
Reference \texttt{.compmodel} files were authored manually from each paper and serve
as the gold standard. Gold models range from simple 4-compartment SIR-like structures
(cholera, flu) to complex multi-group models (HIV with three risk groups, measles with
age stratification), spanning 4--23 flows and 6--19 parameters.
All gold models, extracted drafts, and Phase~3 outputs are available in the project
repository.

\paragraph{LLM Backends.}
Phase~2 uses three backends: OpenAI GPT-5.4-mini~\cite{achiam2023gpt4}, Google
Gemini~2.5 Flash~\cite{team2024gemini}, and Anthropic Claude Opus~4~\cite{anthropic2024claude}.
Phase~3 inference (Tier~3) was evaluated independently with all three backends
(Gemini~2.5 Flash, Claude Opus~4, and OpenAI GPT-5.4-mini) to assess whether the
choice of LLM for gap filling affects outcome.
Temperature is set to 0 for reproducibility in all extraction calls.

\paragraph{Evaluation.}
Recall, precision, and F1 are computed using the fuzzy string matching protocol
described in Section~\ref{sec:phase2eval}. The completeness score
(Equation~\ref{eq:completeness}) is computed per disease over the winner Phase~3 mode.

\subsection{Phase 2: LLM Extraction Quality}
\label{sec:phase2results}

Table~\ref{tab:phase2} reports recall and F1 for all three backends across 10 diseases
under the \emph{enhanced} pipeline and fuzzy string matching.
\emph{Recall} is the primary metric: it measures how many gold-standard entities were
successfully captured, directly quantifying the specification--extraction gap.

\begin{table}[t]
\caption{Phase~2 \textbf{enhanced} extraction results (fuzzy string matching).
R = Recall, F1 = F1-score. C = Compartments, Par = Parameters, F = Flows.
Best recall per entity type across backends is \textbf{bolded}.}
\label{tab:phase2}
\small
\setlength{\tabcolsep}{3.5pt}
\begin{tabular}{llccccccccc}
\toprule
 & & \multicolumn{3}{c}{OpenAI} & \multicolumn{3}{c}{Gemini} & \multicolumn{3}{c}{Claude} \\
\cmidrule(lr){3-5}\cmidrule(lr){6-8}\cmidrule(lr){9-11}
Disease & & C & Par & F & C & Par & F & C & Par & F \\
\midrule
Cholera    & R  & \textbf{1.00} & \textbf{0.83} & \textbf{1.00} & \textbf{1.00} & \textbf{0.83} & \textbf{1.00} & \textbf{1.00} & \textbf{0.83} & \textbf{1.00} \\
           & F1 & 1.00 & 0.91 & 0.86 & 1.00 & 0.91 & 1.00 & 1.00 & 0.91 & 0.60 \\
COVID-19   & R  & \textbf{1.00} & 0.00 & \textbf{1.00} & 0.88 & 0.00 & 0.82 & \textbf{1.00} & 0.00 & \textbf{1.00} \\
           & F1 & 0.70 & 0.00 & 0.96 & 0.61 & 0.00 & 0.86 & 0.70 & 0.00 & 1.00 \\
Dengue     & R  & \textbf{1.00} & \textbf{0.62} & \textbf{1.00} & \textbf{1.00} & \textbf{0.62} & \textbf{1.00} & 0.86 & \textbf{0.62} & 0.86 \\
           & F1 & 1.00 & 0.70 & 0.64 & 1.00 & 0.70 & 1.00 & 0.92 & 0.70 & 0.63 \\
Ebola      & R  & \textbf{1.00} & \textbf{0.07} & \textbf{0.88} & \textbf{1.00} & \textbf{0.07} & \textbf{0.88} & \textbf{1.00} & \textbf{0.07} & 0.50 \\
           & F1 & 1.00 & 0.07 & 0.82 & 1.00 & 0.07 & 0.93 & 1.00 & 0.07 & 0.67 \\
Flu        & R  & \textbf{1.00} & \textbf{1.00} & 0.40 & \textbf{1.00} & \textbf{1.00} & 0.60 & \textbf{1.00} & \textbf{1.00} & \textbf{0.80} \\
           & F1 & 1.00 & 0.38 & 0.57 & 0.80 & 0.43 & 0.75 & 0.80 & 0.32 & 0.89 \\
HIV        & R  & 0.78 & \textbf{0.29} & 0.60 & \textbf{0.89} & \textbf{0.29} & \textbf{0.70} & 0.78 & 0.21 & 0.60 \\
           & F1 & 0.82 & 0.28 & 0.71 & 0.94 & 0.28 & 0.82 & 0.82 & 0.20 & 0.60 \\
Malaria    & R  & \textbf{1.00} & \textbf{0.53} & \textbf{1.00} & \textbf{1.00} & 0.47 & \textbf{1.00} & \textbf{1.00} & 0.47 & \textbf{1.00} \\
           & F1 & 1.00 & 0.48 & 0.92 & 1.00 & 0.47 & 0.96 & 1.00 & 0.47 & 0.96 \\
Measles    & R  & \textbf{1.00} & \textbf{0.42} & 0.85 & \textbf{1.00} & 0.26 & \textbf{1.00} & 0.92 & 0.26 & \textbf{1.00} \\
           & F1 & 1.00 & 0.46 & 0.92 & 1.00 & 0.29 & 1.00 & 0.92 & 0.28 & 0.98 \\
TB         & R  & \textbf{1.00} & \textbf{0.44} & \textbf{0.50} & \textbf{1.00} & \textbf{0.44} & \textbf{0.50} & \textbf{1.00} & \textbf{0.44} & \textbf{0.50} \\
           & F1 & 1.00 & 0.33 & 0.67 & 1.00 & 0.47 & 0.67 & 1.00 & 0.47 & 0.67 \\
Zika       & R  & \textbf{1.00} & 0.26 & \textbf{1.00} & \textbf{1.00} & 0.42 & \textbf{1.00} & \textbf{1.00} & \textbf{0.47} & \textbf{1.00} \\
           & F1 & 1.00 & 0.26 & 1.00 & 1.00 & 0.43 & 1.00 & 1.00 & 0.47 & 1.00 \\
\midrule
\textbf{Avg R} & & \textbf{0.98} & \textbf{0.45} & 0.82 & \textbf{0.98} & 0.44 & \textbf{0.85} & 0.96 & 0.44 & 0.83 \\
\bottomrule
\end{tabular}
\end{table}

Several patterns emerge from the enhanced pipeline results:

\noindent\textbf{Compartments are consistently well-extracted.}
All three backends achieve compartment recall $\geq 0.96$ on average.
OpenAI (GPT-5.4-mini) and Gemini both reach 0.98, with Claude at 0.96.
Perfect compartment recall is achieved for 9 of 10 diseases by OpenAI and Gemini.
Structurally, compartments are the easiest entity type because their names appear
explicitly as state variables in the methods section.

\noindent\textbf{Parameters remain the hardest entity type.}
Even in the enhanced pipeline, parameter recall averages only 0.45 (OpenAI), 0.44
(Gemini), and 0.44 (Claude). COVID-19 shows zero parameter recall across all backends
because that paper embeds parameter values inside ODE derivations rather than naming
them as separate constants. Ebola shows near-zero recall (0.07) across all backends
because parameters are implicit in the derivation. Malaria parameter recall is capped
at 0.53 (OpenAI) because the paper defines many rates as composite symbolic expressions.

\noindent\textbf{Flow recall is competitive across all backends.}
Gemini reaches 0.85 average flow recall, Claude 0.83, and OpenAI 0.82 with GPT-5.4-mini.
The most notable improvement over the previous OpenAI backend (GPT-4o-mini, avg.\ 0.59)
is for malaria (0.09 $\to$ 1.00) and measles (0.00 $\to$ 0.85), diseases where the
stronger model correctly recovers multi-group transition paths implicit in the text.

\noindent\textbf{Provider selection matters.}
Table~\ref{tab:best_provider} shows the winning provider per disease (best combined
score). Claude wins 6 of 10 diseases, OpenAI wins 2 (Malaria, Measles), and Gemini wins 2
(Influenza, HIV).

\paragraph{Original vs.\ enhanced pipeline.}
Table~\ref{tab:phase2} shows results for the enhanced pipeline.
The original pipeline (PyPDF2, separate LLM calls) achieved mean recall of
0.88 / 0.42 / 0.73 (Gemini C/P/F) and 0.92 / 0.47 / 0.78 (Claude).
The enhanced pipeline raises compartment recall by $+10\%$ for Gemini and $+4\%$
for Claude, and flow recall by $+12\%$ for Gemini and $+5\%$ for Claude.
Disease-specific gains include: Ebola flow recall (Gemini) rising from 0.50 to 0.88,
and Malaria compartment recall (Gemini) from 0.29 to 1.00.
For the OpenAI backend, both the pipeline upgrade and the model change
(GPT-4o-mini $\to$ GPT-5.4-mini) contribute jointly to gains, particularly for malaria and
measles flows. Parameter recall improves modestly ($+0.02$--$+0.08$) because the failure
mode (values buried in equations) is not addressed by layout or extraction changes.
All Phase~3 runs use the enhanced pipeline's best-per-disease draft as input.

\begin{table}[t]
\caption{Best Phase~2 provider per disease (enhanced pipeline, composite score:
mean of compartment and flow F1). Score $= \frac{1}{2}(\text{cF1}+\text{fF1})\times 100$.}
\label{tab:best_provider}
\small
\begin{tabular}{llr}
\toprule
Disease & Best Provider & Phase~2 Score \\
\midrule
Cholera         & Claude  & 98.2  \\
COVID-19        & Claude  & 80.0  \\
Dengue          & Claude  & 93.9  \\
Ebola           & Claude  & 81.4  \\
Influenza       & Gemini  & 88.6  \\
HIV             & Gemini  & 85.5  \\
Malaria         & OpenAI  & 89.5  \\
Measles         & OpenAI  & 89.1  \\
Tuberculosis    & Claude  & 89.4  \\
Zika            & Claude  & 89.5  \\
\bottomrule
\end{tabular}
\end{table}

\subsection{Phase 3: Gap Reduction and Recall Improvement}
\label{sec:phase3results}

\paragraph{Ablation study: Phase 2 baseline vs.\ Phase 3 modes.}
In our evaluation, the best Phase~2 draft serves as the baseline extraction system,
representing a pure LLM-based paper-to-model transformation without any post-processing.
Phase~3 then measures the gain from model-driven completion, validation, and repair
over that baseline using paper-grounded retrieval and structural rules.

Table~\ref{tab:ablation} compares the Phase~2 best draft against all three Phase~3
modes using \emph{recall} as the primary metric---the fraction of gold-standard
entities that appear in the extracted model.
Recall is the right metric here because we care about \emph{coverage}: not missing
a required model element. Precision is secondary: an extracted model that contains
extra entities (some of which may be LLM-inferred) is not wrong---it may simply be
more complete than our reference model, which itself is only one possible
interpretation of each paper.
Each cell shows remaining gaps~$g$, compartment recall~$R_C$, and flow recall~$R_F$.
The winner per disease (highest $\bar{R} = (R_C+R_F)/2$) is marked $\star$.
All four modes are evaluated against the same gold standard using the same fuzzy
matching protocol on \texttt{model\_repaired.compmodel}.

\begin{table*}[t]
\caption{Ablation: Phase~2 best draft vs.\ three Phase~3 modes, evaluated by
\textbf{recall} against the gold-standard reference.
$g$ = remaining gaps; $R_C$ = compartment recall; $R_F$ = flow recall;
$\bar{R}$ = mean recall. $\star$ = winner per disease (highest $\bar{R}$).}
\label{tab:ablation}
\small
\setlength{\tabcolsep}{4pt}
\begin{tabular}{lrcccrcccrcccrcc}
\toprule
 & \multicolumn{3}{c}{\textbf{Phase 2}} & \multicolumn{4}{c}{\textbf{P3: RAG-only}} & \multicolumn{4}{c}{\textbf{P3: LLM-only}} & \multicolumn{4}{c}{\textbf{P3: Both}} \\
\cmidrule(lr){2-4}\cmidrule(lr){5-8}\cmidrule(lr){9-12}\cmidrule(lr){13-16}
Disease & $g$ & $R_C$ & $R_F$ & $g$ & $R_C$ & $R_F$ & $\bar{R}$ & $g$ & $R_C$ & $R_F$ & $\bar{R}$ & $g$ & $R_C$ & $R_F$ & $\bar{R}$ \\
\midrule
Cholera      & 4  & 1.00 & 1.00 & 2 & 1.00 & 1.00 & \textbf{1.000}$^\star$ & 2 & 1.00 & 1.00 & \textbf{1.000}$^\star$ & 2 & 1.00 & 1.00 & \textbf{1.000}$^\star$ \\
COVID-19     & 4  & 0.93 & 0.95 & 2 & 1.00 & 1.00 & \textbf{1.000}$^\star$ & 2 & 0.93 & 0.95 & 0.940 & 2 & 0.93 & 0.95 & 0.940 \\
Dengue       & 8  & 1.00 & 0.75 & 4 & 1.00 & 1.00 & \textbf{1.000}$^\star$ & 4 & 1.00 & 1.00 & \textbf{1.000}$^\star$ & 4 & 1.00 & 1.00 & \textbf{1.000}$^\star$ \\
Ebola        & 28 & 0.67 & 0.30 & 19& 0.83 & 0.70 & \textbf{0.767}$^\star$ & 19& 0.83 & 0.70 & \textbf{0.767}$^\star$ & 19& 0.83 & 0.70 & \textbf{0.767}$^\star$ \\
Flu          & 1  & 1.00 & 0.83 & 1 & 1.00 & 1.00 & \textbf{1.000}$^\star$ & 1 & 1.00 & 1.00 & \textbf{1.000}$^\star$ & 1 & 1.00 & 1.00 & \textbf{1.000}$^\star$ \\
HIV          & 18 & 0.78 & 0.38 & 9 & 1.00 & 0.81 & \textbf{0.906}$^\star$ & 9 & 1.00 & 0.81 & \textbf{0.906}$^\star$ & 9 & 0.89 & 0.81 & 0.851 \\
Malaria      & 15 & 1.00 & 1.00 & 12& 1.00 & 1.00 & \textbf{1.000}$^\star$ & 12& 1.00 & 1.00 & \textbf{1.000}$^\star$ & 12& 1.00 & 1.00 & \textbf{1.000}$^\star$ \\
Measles      & 16 & 0.92 & 0.92 & 8 & 1.00 & 1.00 & \textbf{1.000}$^\star$ & 8 & 0.92 & 0.92 & 0.917 & 8 & 0.92 & 0.92 & 0.917 \\
TB           & 5  & 1.00 & 0.50 & 4 & 1.00 & 0.67 & \textbf{0.833}$^\star$ & 4 & 1.00 & 0.67 & \textbf{0.833}$^\star$ & 4 & 1.00 & 0.67 & \textbf{0.833}$^\star$ \\
Zika         & 26 & 1.00 & 1.00 & 13& 1.00 & 1.00 & \textbf{1.000}$^\star$ & 13& 1.00 & 1.00 & \textbf{1.000}$^\star$ & 13& 1.00 & 1.00 & \textbf{1.000}$^\star$ \\
\midrule
\textbf{Mean}& 12.5 & 0.93 & 0.76 & 7.4 & \textbf{0.98} & \textbf{0.92} & \textbf{0.951} & 7.4 & 0.97 & 0.90 & 0.936 & 7.4 & 0.96 & 0.90 & 0.931 \\
\bottomrule
\end{tabular}
\end{table*}

Several observations follow from Table~\ref{tab:ablation}:

\noindent\textbf{All three Phase~3 modes substantially improve recall over Phase~2.}
Mean recall rises from 0.846 (Phase~2) to 0.951 / 0.936 / 0.931 for RAG-only,
LLM-only, and Both respectively---gains of $+10.5\%$, $+9.0\%$, and $+8.5\%$.
All three modes cut mean gap count nearly in half (12.5 $\to$ 7.4).

\noindent\textbf{LLM inference is competitive with RAG on recall.}
When recall is the metric, LLM-only matches RAG-only for 8 of 10 diseases.
It falls short only for COVID-19 ($R_C$: 0.93 vs 1.00) and Measles ($\bar{R}$: 0.917
vs 1.000), where RAG successfully retrieves textually explicit compartments that LLM
inference misses. The gap is small (mean 0.936 vs 0.951), confirming that LLM
inference does add structurally useful entities---it just also adds extras that
lower precision, which recall does not penalize.

\noindent\textbf{The Both mode is rarely strictly better than RAG-only.}
Combining RAG and LLM adds no recall benefit over RAG-only in any disease and slightly
hurts recall for HIV (0.851 vs 0.906 RAG-only): combining modes can displace a
correctly RAG-placed compartment. RAG-only is therefore the safest default.

\noindent\textbf{RAG-only recall is LLM-backend agnostic.}
Phase~3 was re-run with all three LLM backends (Gemini~2.5 Flash, Claude Opus~4,
OpenAI GPT-5.4-mini) for both Phase~3 inference modes.
The \emph{rag\_only} winner mode produces \emph{identical} recall values across all
three backends in every disease: the same winner mode is selected (rag\_only for 9 of 10,
both for Malaria) and the same $\Delta$Recall is achieved ($+7.6\% / +7.6\% / +20.8\%$
for compartments, parameters, and flows respectively).
Only the \emph{llm\_only} mode shows backend sensitivity: OpenAI produces lower F1 on
Cholera (cF1~0.67 vs 0.89 for Gemini and Claude), reflecting that unconstrained LLM
inference varies by model, while RAG-grounded retrieval does not.
This confirms that the recall improvement in Phase~3 is driven entirely by the
retrieval mechanism, not by the choice of LLM backend---making the approach reproducible
across commercial LLM providers.

\noindent\textbf{Phase 3 always improves or matches recall; it never harms it.}
Unlike F1, recall does not penalize extra entities, so Phase~3 never regresses on this
metric. Tuberculosis (Phase~2 $\bar{R}=0.750$, Phase~3 $0.833$) improves, even though
it showed an F1 regression in the F1-based view. The Phase~3 model correctly adds a
missing flow that was in the TB paper; the F1 drop comes from additional entities that
are in the paper but absent from our reference model---entities that a completeness
tool should recover, not suppress.

\noindent\textbf{Ebola remains the hardest disease.}
Even after Phase~3, Ebola reaches only $\bar{R}=0.813$ (compartments 1.0, flows 0.625).
Three of eight Ebola flows remain missing---a direct consequence of the paper's
implicit flow specification discussed in Section~\ref{sec:discussion}.

\paragraph{Phase 2 vs Phase 3 recall improvement.}
Table~\ref{tab:recall_delta} compares Phase~2 and Phase~3 recall (winner mode)
against the same gold standard using identical fuzzy matching. The gap delta column
shows the net number of gaps closed by Phase~3.

\begin{table}[t]
\caption{Phase~2 enhanced best-per-disease draft vs Phase~3 filled recall, compared
against the gold-standard reference using the same fuzzy string matching.
Gap $\Delta$ = gaps closed by Phase~3 (positive = improvement).}
\label{tab:recall_delta}
\small
\setlength{\tabcolsep}{4pt}
\begin{tabular}{lccccccr}
\toprule
 & \multicolumn{3}{c}{Phase~2 Recall (C/P/F)} & \multicolumn{3}{c}{Phase~3 Recall (C/P/F)} & Gap \\
Disease & C & P & F & C & P & F & $\Delta$ \\
\midrule
Cholera      & 0.750 & 1.000 & 0.667 & 1.000 & 1.000 & 1.000 & +2 \\
COVID-19     & 0.875 & 1.000 & 0.909 & 1.000 & 1.000 & 1.000 & +2 \\
Dengue       & 1.000 & 1.000 & 0.857 & 1.000 & 1.000 & 1.000 & +4 \\
Ebola        & 0.833 & 0.500 & 0.375 & 1.000 & 0.571 & 0.625 & +4 \\
Flu          & 1.000 & 1.000 & 0.800 & 1.000 & 1.000 & 1.000 & +0 \\
HIV          & 0.778 & 1.000 & 0.300 & 1.000 & 1.000 & 1.000 & +9 \\
Malaria      & 1.000 & 0.737 & 1.000 & 1.000 & 1.000 & 1.000 & +7 \\
Measles      & 1.000 & 0.947 & 0.800 & 1.000 & 1.000 & 1.000 & +6 \\
TB           & 1.000 & 0.889 & 0.500 & 1.000 & 1.000 & 0.667 & +1 \\
Zika         & 1.000 & 0.737 & 1.000 & 1.000 & 1.000 & 1.000 & +13\\
\midrule
\textbf{Mean $\Delta$} & \multicolumn{3}{c}{} & \multicolumn{3}{c}{+0.076 / +0.076 / +0.208} & \\
\bottomrule
\end{tabular}
\end{table}

Phase~3 achieves perfect recall ($= 1.0$) for compartments and flows in 9 of 10 diseases
and for all three entity types in 7 of 10 diseases.
Phase~3's dominant contribution is now flow recovery: mean flow recall rises
from $0.721$ (Phase~2 draft) to $0.929$ (Phase~3), a gain of $+20.8\%$.
Compartment recall improves by $+7.6\%$ (from $0.924$ to $1.000$ in most diseases)
and parameter recall improves by $+7.6\%$ (from $0.881$ to $0.957$ mean).
The shift compared to earlier pipeline versions reflects that with the stronger
GPT-5.4-mini backend the best-per-disease Phase~2 draft already captures many
parameters directly, leaving flow and compartment gaps as the primary Phase~3 targets.

The one disease where Phase~3 adds no flow improvement (\emph{flu}) already has
all flows filled by the Phase~2 draft; the remaining gap is a parameter value that RAG
cannot locate because the flu paper reports it only in a figure caption.
Ebola remains the hardest disease: even after Phase~3, compartment recall reaches 1.0
but flow recall reaches only 0.625, because flow endpoints are implicit in the derivation.

\subsection{Phase 3 Completeness Scores}

Table~\ref{tab:completeness} breaks down the completeness score per component
for the winner Phase~3 run of each disease.

\begin{table}[t]
\caption{Phase~3 completeness scores (winner mode = rag\_only for 9 diseases, both for Malaria).
G = gap reduction \%, R = reference agreement \%, T = fill traceability \%,
A = parameter accuracy \%, S = structural integrity \%. Score = Eq.~\eqref{eq:completeness}.}
\label{tab:completeness}
\small
\setlength{\tabcolsep}{4pt}
\begin{tabular}{lrrrrrrr}
\toprule
Disease & G & R & T & A & S & \textbf{Score} \\
\midrule
Cholera      & 100 & 71.7 & 50.0 & 100 & 50.0 & 75.4 \\
COVID-19     & 100 & 97.2 & 50.0 & 100 & 51.3 & 82.0 \\
Dengue       & 100 & 90.9 & 100  & 100 & 41.7 & 89.0 \\
Ebola        & 33.3& 74.2 & 37.5 & 100 & 40.0 & 55.4 \\
Flu          & 0.0 & 83.9 & 100  & 100 & 47.4 & 63.1 \\
HIV          & 100 & 89.0 & 77.8 & 100 & 25.0 & 81.6 \\
Malaria      & 100 & 95.8 & 28.6 & 28.6& 26.5 & 62.9 \\
Measles      & 100 & 98.8 & 100  & 50.0& 50.0 & 84.7 \\
Tuberculosis & 25.0& 70.8 & 100  & 100 & 54.5 & 67.1 \\
Zika         & 100 & 85.6 & 84.6 & 71.4& 36.7 & 79.5 \\
\midrule
\textbf{Mean}& 75.8& 85.8 & 72.9 & 85.0& 42.3 & \textbf{74.0} \\
\bottomrule
\end{tabular}
\end{table}

\paragraph{Component analysis.}
\textbf{Reference agreement} is the highest-scoring component (mean 85.8\%), confirming
that Phase~3 fills bring model structure close to the gold standard.
\textbf{Gap reduction} improved to mean 75.8\%, with Malaria now achieving 100\% gap
reduction (winner mode = both), compared to 25\% in the previous run where fewer gaps
were identified.
\textbf{Parameter accuracy} averages 85.0\%; the main drag is Malaria (28.6\%), where
Phase~3 fills parameters by name but the paper's composite symbolic expressions do not
yield exact numeric matches against the gold standard.
\textbf{Fill traceability} averages 72.9\%; the main drag is Ebola (37.5\%), where most
Phase~3 fills are partially flagged rather than fully retrieved---indicating the Ebola
paper's parameter specification is substantially implicit.
\textbf{Structural integrity} is the weakest component (mean 42.3\%), because several
model classes systematically resist rule-based repair:
multi-group HIV models have up to 14 parameters that cannot be automatically assigned
to their correct flows after a collapse fix; the Zika and Dengue host--vector models
involve composite expressions that the repairer cannot safely decompose.

\paragraph{Interpreting the 74.0\% mean score.}
The completeness score is not a grade for the extracted models---it is a measurement
instrument for the papers themselves. A paper that scores 100 would be one that fully
specifies every compartment, parameter, and flow with numeric values traceable to its
text and a structurally consistent model. No paper in our dataset achieves this.
A score of 74.0\% means that three-quarters of the gold-standard structure is
recoverable from the paper text under systematic extraction---and that the remaining
quarter represents genuine specification gaps in the published literature.

Before this framework, that number did not exist: there was no instrument to compute
it. Researchers either trusted published models as complete or noted incompleteness
anecdotally. Our contribution is to make the gap measurable so it can be studied,
compared across diseases and papers, and used to argue for more explicit specification
practices.

A model can also achieve high recall by name-matching while still having orphaned
parameters, incorrect flow rates, or untraced fills. Completeness is intentionally
stricter: it rewards coverage, traceability, accuracy, and structural soundness jointly,
so that a model which looks correct on the surface but is internally inconsistent is
penalized.

\subsection{Structural Error Analysis}

Across all 10 Phase~3 models, the structural validator detected 192 errors before repair.
After the repair pass, 107 errors remain---a resolution rate of 44\%.
Table~\ref{tab:structural} summarizes error counts and resolution by type.

\begin{table}[t]
\caption{Structural error counts before and after Phase~3 repair, aggregated over 10 diseases.}
\label{tab:structural}
\small
\begin{tabular}{lrrr}
\toprule
Error Type & Before & After & Resolved (\%) \\
\midrule
Self-referential flow       & 12 & 9  & 25\% \\
Orphaned parameters         & 88 & 60 & 32\% \\
Uniform param.\ collapse    & 10 & 0  & 100\% \\
Zero population             & 10 & 0  & 100\% \\
Broken flow chain           & 18 & 10 & 44\% \\
Missing birth source        & 30 & 12 & 60\% \\
Parameter contamination     & 24 & 16 & 33\% \\
\midrule
\textbf{Total}              & 192& 107& \textbf{44\%} \\
\bottomrule
\end{tabular}
\end{table}

Deterministic errors (\emph{uniform parameter collapse}, \emph{zero population}) are
fully resolved because their fixes are unconditional. Orphaned parameters are the most
prevalent class and the hardest to resolve: multi-group and vector models declare
many group-specific rate parameters (e.g., $b_{hw}$, $c_{hm}$ in HIV) that cannot be
automatically wired when all flow slots are already occupied.
Self-referential flows are annotated rather than structurally fixed when no
intermediate (Exposed/Latent) compartment is available to redirect the flow through.
